\subsubsection{A common seven-unit theta heat band from the sine majorant}
\label{subsec:sevenunit}
The quadratic cosine minorant leaves a positive band
$|x|<\sqrt{48}$. The same previously proved variance
estimate admits a strictly stronger consequence.
Write the Fourier-correlation representation as
\[
U_n(s,x)=\int_{\mathbb R}K_{n,s}(w)\cos(xw)\,dw,
\qquad K_{n,s}(w)\ge0,
\]
with $A_n(s)=\int K_{n,s}>0$ and
$B_n(s)=\int w^2K_{n,s}$.
By the established theta moment estimate,
\begin{equation}\label{eq:sevenvar}
0\le B_n(s)/A_n(s)\le1/24
\quad(n\ge1,\ 0\le s\le1/2).
\end{equation}

\begin{theorem}[Seven-unit division-free heat comparison]
\label{thm:sevenunit}
For the genuine theta heat family, every $n\ge1$
and $0\le s\le1/2$ satisfy, with $r=|x|$,
\begin{equation}\label{eq:sevenpiece}
\frac{U_n(s,x)}{A_n(s)}
\ge g(r):=
\begin{cases}
1-r^2/48,&0\le r\le\sqrt{24},\\[2pt]
\frac32-r/\sqrt{24},&r\ge\sqrt{24}.
\end{cases}
\end{equation}
Consequently $U_n(s,x)>0$ for
$|x|<3\sqrt6$.
In particular, for the whole \emph{closed}
seven-unit band $|x|\le7$,
\[
\boxed{\quad U_n(s,x)\ge A_n(s)/15>0,\qquad
k_n U_{n+1}(s,x)\le
15k_n\frac{A_{n+1}(s)}{A_n(s)}U_n(s,x).\quad}
\]
This comparison is division-free with respect
to noncentral Laguerre values.
\end{theorem}
\begin{proof}
Differentiating under the absolutely convergent
correlation integral and using both
$|\sin v|\le |v|$ and $|\sin v|\le1$ gives
\[
|U_n'(s,x)|\le
\min\!\left(rB_n(s),\int|w|K_{n,s}(w)\,dw\right).
\]
By Cauchy--Schwarz and \eqref{eq:sevenvar},
\[
|U_n'(s,x)|\le A_n(s)
\min\!\left(\frac r{24},\frac1{\sqrt{24}}\right).
\]
Integrating this derivative bound from $0$
to $r$ gives \eqref{eq:sevenpiece}.
The second branch is positive precisely when
$r<(3/2)\sqrt{24}=3\sqrt6$.

For the rational seven-unit statement,
$7<3\sqrt6$ and
\[
\left(\frac{43}{30}\right)^2\,24-49
=\frac{23}{75}>0.
\]
Since both sides compared are positive,
this implies
$7/\sqrt{24}<43/30$ and therefore
$g(7)>1/15$.
Monotonicity of $g$ yields
$U_n(s,x)\ge A_n(s)/15$ throughout
$|x|\le7$. The positive correlation
also gives $U_{n+1}(s,x)\le A_{n+1}(s)$.
Multiplying the lower bound for $U_n$
by $15k_nA_{n+1}/A_n$ gives the
claimed inequality.
\end{proof}
The interval $|x|\le7$ strictly improves the
six-unit and $\sqrt{48}$ bands using no
new theta integration. It remains a finite
region; the argument is not a global RH proof.
The relevant rational comparison gates are
recorded in
\path{XiKernel/SevenUnitHeatBand.lean}
and \path{certificates/xi_seven_unit_heat_exact.py}.
The analytic source inputs are the positive
correlation measure and the already established
theta moment estimate, not an end-to-end Lean
proof of the theta integrals.
